\documentclass[12pt,a4paper]{article}
\usepackage[utf8]{inputenc}
\usepackage[T1]{fontenc}
\usepackage[french]{babel}
\usepackage{amsmath}
\usepackage{amsfonts}
\usepackage{amssymb}
\usepackage{graphicx}
\usepackage{xcolor}
\usepackage{tcolorbox}
\usepackage{enumitem}
\usepackage{pifont}
\usepackage{fancyhdr}
\usepackage{geometry}
\usepackage{listings}
\usepackage{hyperref}
\usepackage{float}
\usepackage{charter}
\usepackage{csquotes}
\usepackage{tikz}
\usetikzlibrary{shapes,positioning,shadows,decorations.pathmorphing}

% Configuration de la géométrie
\geometry{left=2cm,right=2cm,top=2.5cm,bottom=2.5cm}

% Configuration des couleurs
\definecolor{codegreen}{rgb}{0,0.6,0}
\definecolor{codegray}{rgb}{0.5,0.5,0.5}
\definecolor{codepurple}{rgb}{0.58,0,0.82}
\definecolor{backcolour}{rgb}{0.95,0.95,0.92}
\definecolor{warningred}{RGB}{204,0,0}
\definecolor{successgreen}{RGB}{0,153,0}
\definecolor{pyblue}{RGB}{55,118,171}
\definecolor{pyyellow}{RGB}{255,212,59}
\definecolor{spyderred}{RGB}{200,50,50}

\lstset{
	backgroundcolor=\color{backcolour},
	commentstyle=\color{codegreen},
	keywordstyle=\color{red},
	numberstyle=\tiny\color{codegray},
	stringstyle=\color{codepurple},
	basicstyle=\footnotesize\ttfamily,
	breakatwhitespace=false,
	breaklines=true,
	captionpos=b,
	keepspaces=true,
	numbers=left,
	numbersep=5pt,
	showspaces=false,
	showstringspaces=false,
	showtabs=false,
	tabsize=2,
	frame=single,
	framerule=1pt,
	rulecolor=\color{codegray}
}

% Boîtes personnalisées avec bords arrondis
\newtcolorbox{astuce}{
	colback=yellow!10!white,
	colframe=yellow!50!black,
	title=Astuce,
	fonttitle=\bfseries,
	sharp corners,
	rounded corners=all=5mm
}

\newtcolorbox{attention}{
	colback=red!5!white,
	colframe=red!50!black,
	title=Attention~!,
	fonttitle=\bfseries,
	sharp corners,
	rounded corners=all=5mm
}

\newtcolorbox{reussite}{
	colback=green!5!white,
	colframe=green!50!black,
	title=Bravo~!,
	fonttitle=\bfseries,
	sharp corners,
	rounded corners=all=5mm
}

\newtcolorbox{definition}{
	colback=blue!5!white,
	colframe=blue!50!black,
	title=Définition,
	fonttitle=\bfseries,
	sharp corners,
	rounded corners=all=5mm
}

\newtcolorbox{etape}{
	colback=white,
	colframe=orange!50!black,
	title=\textbf{Étape},
	fonttitle=\bfseries,
	sharp corners,
	rounded corners=all=5mm
}

% Configuration des en-têtes
\pagestyle{fancy}
\fancyhf{}
\fancyhead[L]{\leftmark}
\fancyhead[R]{\thepage}
%\fancyfoot[C]{\copyright~Abdallah Houmadi SOIRFANE - \the\year}

% ------------------------------------------------------------
% Logos stylisés dessinés en TikZ (aucune dépendance réseau).
% Remplacent les images distantes de l'original, que LaTeX ne
% peut de toute façon pas charger depuis une URL directement.
% ------------------------------------------------------------
\newcommand{\logopython}{%
	\begin{tikzpicture}
		\node[circle,fill=pyblue,minimum size=2.6cm,text=white,font=\bfseries\Large,align=center] {Py};
	\end{tikzpicture}%
}
\newcommand{\logospyder}{%
	\begin{tikzpicture}
		\node[circle,fill=spyderred,minimum size=2.6cm,text=white,font=\bfseries\Large,align=center] {Spy};
	\end{tikzpicture}%
}

% Commande pour les captures d'écran locales (à fournir par
% l'utilisateur dans le dossier img/ -- voir annexe technique)
\newcommand{\capture}[3]{%
	\begin{figure}[H]
		\centering
		\includegraphics[width=0.8\textwidth]{#1}
		\caption{#2}
		\label{fig:#3}
	\end{figure}
}

\title{\Huge\textbf{Guide d'Installation Python et Spyder}}
\author{\Large\textbf{\textit{Abdallah Houmadi SOIRFANE}}}
\date{\Large \today}

\begin{document}
	
	% Page de titre
	\maketitle
	
	
	% ============================================================
	% PROLOGUE
	% ============================================================
	\section*{Prologue~: le problème du jour}
	\addcontentsline{toc}{section}{Prologue~: le problème du jour}
	
	Youssouf a terminé son cours de Python. Il est chaud. Il est prêt. Il veut installer Python et Spyder pour commencer à coder.
	
	Il ouvre Google. Il tape «~installer Python~». Il clique sur le premier lien.
	
	\textbf{DEUX HEURES PLUS TARD...}
	
	Son ordinateur ressemble à un champ de bataille~:
	\begin{itemize}
		\item 3 versions de Python qui se battent dans son PATH
		\item Des messages d'erreur rouges partout
		\item Un dossier «~Téléchargements~» qui contient 47 fichiers .exe inutiles
		\item Et une envie irrésistible de jeter son PC par la fenêtre
	\end{itemize}
	
	Il aurait dû lire ce guide. Mais il ne l'a pas fait. Ne faites pas comme Youssouf.
	
	\begin{attention}
		\textbf{Ne faites pas comme Youssouf~!} Lisez ce guide en entier avant de commencer l'installation. Il ne vous faudra que 15 minutes, et ça vous évitera deux heures de galère.
	\end{attention}
	
	% ============================================================
	% SECTION 0 : VOCABULAIRE DE BASE
	% ============================================================
	\section{Avant de commencer~: le vocabulaire de base}
	
	Si vous n'avez jamais installé d'outil de développement, certains mots vont revenir sans cesse dans ce guide. Mieux vaut les comprendre \emph{avant} de foncer.
	
	\begin{definition}
		\textbf{Un terminal} (ou «~invite de commandes~» sur Windows) est une fenêtre où vous tapez du texte au lieu de cliquer sur des icônes. C'est l'outil que l'on utilise pour vérifier qu'un programme est bien installé, ou pour en installer un autre. Ce n'est pas réservé aux experts~: on y tape une ligne, on appuie sur Entrée, et l'ordinateur répond.
		
		\textbf{Comment l'ouvrir~?}
		\begin{itemize}
			\item \textbf{Windows} : touche \texttt{Windows}, tapez \texttt{cmd} ou \texttt{PowerShell}, appuyez sur Entrée.
			\item \textbf{Mac} : \texttt{Cmd+Espace}, tapez \texttt{Terminal}, appuyez sur Entrée.
			\item \textbf{Linux} : \texttt{Ctrl+Alt+T} dans la plupart des distributions.
		\end{itemize}
	\end{definition}
	
	\begin{definition}
		\textbf{Le PATH} est une liste d'endroits où votre ordinateur va chercher les programmes quand vous tapez leur nom dans un terminal. Si Python n'est pas dans cette liste, taper \texttt{python} ne fonctionnera pas, même si Python est bien installé sur le disque -- l'ordinateur ne sait simplement pas où le trouver.
	\end{definition}
	
	\begin{definition}
		\textbf{pip} est l'outil qui installe des «~packages~» (des bibliothèques de code écrites par d'autres) pour Python. Spyder, Pandas, Matplotlib sont des packages que l'on installe avec pip. On peut le voir comme le magasin d'applications de Python.
	\end{definition}
	
	\begin{definition}
		\textbf{Python} est un langage de programmation~: le cerveau qui va exécuter vos instructions.
		
		\textbf{Spyder} est un environnement de développement (IDE)~: l'endroit où vous allez écrire votre code, le tester et voir les résultats. Python peut fonctionner sans Spyder, mais Spyder ne sert à rien sans Python~: c'est pourquoi on installe toujours Python en premier.
	\end{definition}
	
	\begin{center}
		\begin{tikzpicture}
			\node[align=center] (py) at (0,0) {\logopython\\[2pt]\textbf{Python}\\[-2pt]\footnotesize le langage};
			\node[align=center] (spy) at (5,0) {\logospyder\\[2pt]\textbf{Spyder}\\[-2pt]\footnotesize l'atelier de travail};
			\draw[->,thick] (py) -- node[above,font=\footnotesize]{est utilisé dans} (spy);
		\end{tikzpicture}
	\end{center}
	
	\begin{astuce}
		Pourquoi Spyder et pas un autre~? Parce que Spyder est spécialement conçu pour la data science et le calcul scientifique. Il a tout ce dont vous avez besoin sans rien de superflu. Pour plus d'informations~: \url{https://www.spyder-ide.org/}
	\end{astuce}
	
	% ============================================================
	% SECTION 2 : INSTALLATION DE PYTHON
	% ============================================================
	\section{Installation de Python}
	
	\subsection{Étape 1~: télécharger Python}
	
	\begin{etape}
		\begin{enumerate}
			\item Ouvrez votre navigateur.
			\item Allez sur~: \texttt{https://www.python.org/downloads/}
			\item Le site détecte automatiquement votre système d'exploitation et affiche un gros bouton \textbf{Download Python}.
			\item Cliquez dessus pour lancer le téléchargement.
		\end{enumerate}
	\end{etape}
	
	\begin{astuce}
		Python est gratuit. Si on vous demande de payer, vous n'êtes pas sur le bon site. L'URL officielle est \texttt{python.org} -- méfiez-vous des sites qui ressemblent à s'y méprendre mais dont l'adresse est légèrement différente.
	\end{astuce}
	
	\subsection{Étape 2~: installer Python}
	
	\subsubsection{Sur Windows}
	
	\begin{etape}
		\begin{enumerate}
			\item Double-cliquez sur le fichier \texttt{python-3.x.x.exe} téléchargé.
			
			\item \textbf{LE MOMENT CRUCIAL~:} en bas de la première fenêtre, cochez \textbf{\textcolor{red}{« Add python.exe to PATH »}}.
			
			\item Cliquez sur \textbf{Install Now}.
			\item L'installation commence (elle dure une à deux minutes).
			\item À la fin, cliquez sur \textbf{Close}.
		\end{enumerate}
	\end{etape}
	
	\begin{center}
		\begin{tikzpicture}[scale=1]
			\draw[thick,rounded corners,fill=gray!5] (0,0) rectangle (8,3.4);
			\draw[fill=gray!20] (0,3.0) rectangle (8,3.4);
			\node[font=\footnotesize\bfseries] at (4,3.2) {Python 3.x.x Setup};
			\node[font=\footnotesize,align=left] at (4,2.4) {Install Python 3.x.x (64-bit)};
			\draw (0.6,1.55) rectangle (1.0,1.95);
			\draw[thick] (0.65,1.6) -- (0.8,1.45+0.05) -- (1.0,1.95);
			\node[font=\footnotesize,anchor=west,text=red,align=left] at (1.15,1.75) {\bfseries Add python.exe to PATH};
			\draw[fill=pyblue,rounded corners] (2.0,0.4) rectangle (4.2,0.9);
			\node[font=\footnotesize,text=white] at (3.1,0.65) {Install Now};
			\draw[rounded corners] (4.5,0.4) rectangle (6.2,0.9);
			\node[font=\footnotesize] at (5.35,0.65) {Customize};
		\end{tikzpicture}
	\end{center}
	
	\begin{attention}
		\textbf{Pourquoi cocher « Add python.exe to PATH »~?}
		
		C'est LE piège classique. Sans cette case, votre ordinateur ne saura pas où trouver Python quand vous tapez \texttt{python} dans le terminal, même si l'installation s'est bien passée. C'est comme avoir une voiture flambant neuve, mais garée dans un garage dont vous avez perdu l'adresse. Si vous avez oublié de cocher la case, pas de panique~: relancez simplement l'installateur, choisissez «~Modify~», puis cochez la case (voir la section Dépannage).
	\end{attention}
	
	\subsubsection{Sur Mac}
	
	\begin{etape}
		\begin{enumerate}
			\item Double-cliquez sur le fichier \texttt{python-3.x.x-macosx.pkg} téléchargé.
			\item Cliquez sur \textbf{Continuer} et acceptez la licence.
			\item Cliquez sur \textbf{Installer} et entrez votre mot de passe administrateur.
			\item Cliquez sur \textbf{Fermer}.
		\end{enumerate}
		Sur Mac, Python est automatiquement ajouté au PATH~: pas de case à cocher.
	\end{etape}
	
	\subsubsection{Sur Linux}
	
	\begin{etape}
		Python est très souvent déjà installé. Vérifiez d'abord avec~:
		\begin{lstlisting}[language=bash]
			python3 --version
		\end{lstlisting}
		Si la commande n'est pas reconnue, installez-le avec~:
		\begin{lstlisting}[language=bash]
			sudo apt update
			sudo apt install python3 python3-pip
		\end{lstlisting}
		(Remplacez \texttt{apt} par \texttt{dnf} sur Fedora, ou \texttt{pacman -S} sur Arch, selon votre distribution.)
	\end{etape}
	
	\subsection{Étape 3~: vérifier l'installation}
	
	\begin{etape}
		\begin{enumerate}
			\item Ouvrez un terminal (voir section précédente si vous ne savez pas comment faire).
			\item Tapez~:
			\begin{lstlisting}[language=bash]
				python --version
			\end{lstlisting}
			\item Vous devriez voir~: \texttt{Python 3.x.x}
			\item Tapez aussi~:
			\begin{lstlisting}[language=bash]
				pip --version
			\end{lstlisting}
		\end{enumerate}
		Sur Mac et Linux, si \texttt{python} ne fonctionne pas, essayez \texttt{python3} -- certains systèmes distinguent les deux noms.
	\end{etape}
	
	\begin{reussite}
		\textbf{Vous avez réussi si~:}
		\begin{itemize}
			\item \texttt{python --version} (ou \texttt{python3 --version}) affiche \texttt{Python 3.x.x}
			\item \texttt{pip --version} affiche un numéro de version
		\end{itemize}
		Si l'un des deux échoue, direction la section Dépannage -- ne restez pas bloqué, la solution y est probablement déjà écrite.
	\end{reussite}
	
	% ============================================================
	% SECTION 3 : INSTALLATION DE SPYDER
	% ============================================================
	\section{Installation de Spyder}
	
	Maintenant que Python est installé, on va installer Spyder. Il existe deux méthodes~: choisissez celle qui vous correspond.
	
	\begin{center}
		\begin{tabular}{|p{4.2cm}|p{5.3cm}|p{5.3cm}|}
			\hline
			& \textbf{Méthode 1~: pip} & \textbf{Méthode 2~: Anaconda} \\
			\hline
			Taille du téléchargement & Petite (quelques Mo) & Grande (plusieurs centaines de Mo) \\
			\hline
			Bibliothèques scientifiques (NumPy, Pandas...) & À installer séparément & Déjà incluses \\
			\hline
			Recommandé pour & Ceux qui veulent un système léger et savent déjà ce qu'ils installent & Les débutants qui veulent que «~tout marche du premier coup~» \\
			\hline
		\end{tabular}
	\end{center}
	
	\subsection{Méthode 1~: installation avec pip}
	
	\begin{etape}
		\begin{enumerate}
			\item Ouvrez un terminal.
			\item Tapez~:
			\begin{lstlisting}[language=bash]
				pip install spyder
			\end{lstlisting}
			\item Patientez quelques minutes (Spyder a besoin de plusieurs autres bibliothèques).
			\item Pour lancer Spyder, tapez~:
			\begin{lstlisting}[language=bash]
				spyder
			\end{lstlisting}
		\end{enumerate}
	\end{etape}
	
	\subsection{Méthode 2~: installation via Anaconda (recommandée pour débuter)}
	
	Anaconda est une distribution qui regroupe Python, Spyder et l'essentiel des bibliothèques scientifiques (NumPy, Pandas, Matplotlib) en une seule installation. C'est la méthode la plus fiable pour un débutant.
	
	\begin{etape}
		\begin{enumerate}
			\item Téléchargez Anaconda sur \texttt{https://www.anaconda.com/download}.
			\item Installez-le comme un programme normal (suivant, suivant, installer).
			\item Spyder est déjà installé~: pas besoin de taper la moindre commande.
			\item Lancez-le depuis \textbf{Anaconda Navigator} (une fenêtre avec des icônes), ou directement depuis le menu Démarrer / Launchpad en cherchant «~Spyder~».
		\end{enumerate}
	\end{etape}
	
	\begin{astuce}
		Si vous installez Anaconda, vous n'avez pas besoin de refaire l'installation de Python de la section précédente~: Anaconda l'inclut déjà. Dans ce cas, à l'étape de vérification, tapez simplement \texttt{python --version} dans le terminal \textbf{Anaconda Prompt} (Windows) plutôt que le terminal classique.
	\end{astuce}
	
	% ============================================================
	% SECTION 4 : DÉCOUVERTE DE SPYDER
	% ============================================================
	\section{Découverte de Spyder}
	
	Quand vous ouvrez Spyder pour la première fois, l'écran est découpé en plusieurs zones. Voici à quoi vous attendre~:
	
	\begin{center}
		\begin{tikzpicture}[scale=1,every node/.style={font=\footnotesize}]
			\draw[thick] (0,0) rectangle (10,6);
			% zone editeur (gauche, pleine hauteur)
			\draw[fill=blue!8] (0,0) rectangle (5.2,6);
			\node[align=center] at (2.6,3) {\textbf{1. Éditeur de code}\\(à gauche)\\[4pt]\footnotesize Vous écrivez\\vos programmes ici};
			% zone variables (haut droite)
			\draw[fill=green!10] (5.2,3.2) rectangle (10,6);
			\node[align=center] at (7.6,4.6) {\textbf{2. Explorateur de}\\\textbf{variables}\\(en haut à droite)};
			% zone console (bas droite)
			\draw[fill=orange!10] (5.2,0) rectangle (10,3.2);
			\node[align=center] at (7.6,1.6) {\textbf{3. Console IPython}\\(en bas à droite)\\[4pt]\footnotesize Les résultats\\s'affichent ici};
		\end{tikzpicture}
	\end{center}
	
	\subsection{Les zones principales}
	
	\begin{enumerate}
		\item \textbf{Éditeur de code (à gauche)} : c'est là que vous écrivez vos programmes, ligne par ligne, dans un fichier \texttt{.py}.
		
		\item \textbf{Explorateur de variables (en haut à droite)} : affiche toutes vos variables en cours d'utilisation, même les tableaux Pandas en format tableau -- très utile pour vérifier ce que contient une variable sans taper de code supplémentaire.
		
		\item \textbf{Console IPython (en bas à droite)} : les résultats de votre programme s'affichent ici. Elle est interactive~: vous pouvez aussi y taper des commandes directement, sans passer par un fichier.
		
		\item \textbf{Zone des graphiques} (un onglet à côté de l'explorateur de variables) : affiche les figures générées avec Matplotlib ou d'autres bibliothèques de visualisation.
	\end{enumerate}
	
	\begin{astuce}
		Vous pouvez redimensionner les zones en tirant sur les bordures, et les réorganiser via le menu \textbf{View} $\rightarrow$ \textbf{Panes}. Si vous cassez tout par accident, \textbf{View} $\rightarrow$ \textbf{Window layouts} $\rightarrow$ \textbf{Default layout} remet tout en ordre.
	\end{astuce}
	
	% ============================================================
	% SECTION 5 : VOTRE PREMIER PROGRAMME
	% ============================================================
	\section{Votre premier programme avec Spyder}
	
	\subsection{Étape 1~: créer un fichier}
	
	\begin{etape}
		\begin{enumerate}
			\item Dans Spyder~: \textbf{File} $\rightarrow$ \textbf{New File} (\texttt{Ctrl+N}).
			\item Enregistrez~: \textbf{File} $\rightarrow$ \textbf{Save} (\texttt{Ctrl+S}).
			\item Nommez-le, par exemple \texttt{mon\_premier.py}. L'extension \texttt{.py} indique que c'est un fichier Python.
		\end{enumerate}
	\end{etape}
	
	\subsection{Étape 2~: écrire le code}
	
	\begin{etape}
		Écrivez ceci dans l'éditeur~:
		\begin{lstlisting}[language=Python]
			# Mon premier programme Python dans Spyder
			print("Bonjour le monde !")
			
			nom = "Youssouf"
			age = 17
			total = 3 * 1500 + 2 * 2000
			
			print(f"Total : {total}")
			print(f"Je m'appelle {nom} et j'ai {age} ans")
		\end{lstlisting}
	\end{etape}
	
	\subsection{Étape 3~: exécuter}
	
	\begin{etape}
		\begin{enumerate}
			\item Cliquez sur le bouton \textbf{Run} (la flèche verte) ou appuyez sur \texttt{F5}.
			\item Une fenêtre peut apparaître la première fois pour choisir les paramètres d'exécution~: laissez les valeurs par défaut et cliquez sur \textbf{Run}.
			\item Le résultat s'affiche dans la console, en bas à droite.
		\end{enumerate}
	\end{etape}
	
	\begin{reussite}
		\textbf{FÉLICITATIONS~!} Vous venez d'exécuter votre premier programme Python dans Spyder~!
	\end{reussite}
	
	\subsection{Un deuxième exemple, un peu plus loin}
	
	Pour vous montrer que Python peut aussi réagir à ce que vous tapez et prendre des décisions, voici un second programme à essayer~:
	
	\begin{etape}
		\begin{lstlisting}[language=Python]
			# Un programme qui demande une note et donne une appreciation
			note = float(input("Entrez votre note sur 20 : "))
			
			if note >= 16:
			appreciation = "Excellent"
			elif note >= 12:
			appreciation = "Bien"
			elif note >= 10:
			appreciation = "Passable"
			else:
			appreciation = "Insuffisant"
			
			print(f"Avec {note}/20, votre appreciation est : {appreciation}")
		\end{lstlisting}
		Exécutez-le avec \texttt{F5}. La console va vous demander une valeur~: tapez un nombre, par exemple \texttt{14}, puis appuyez sur Entrée.
	\end{etape}
	
	\begin{astuce}
		Si vous voulez exécuter seulement quelques lignes sans relancer tout le fichier, sélectionnez-les puis appuyez sur \texttt{F9}. C'est très pratique pour tester un petit bout de code isolément.
	\end{astuce}
	
	% ============================================================
	% SECTION 6 : DÉPANNAGE
	% ============================================================
	\section{Dépannage}
	
	\subsection{Erreur~: «~python n'est pas reconnu~»}
	
	\textbf{Cause~:} vous n'avez pas coché «~Add python.exe to PATH~» à l'installation (Windows).
	
	\textbf{Solution~:}
	\begin{enumerate}
		\item Relancez le fichier \texttt{python-3.x.x.exe} téléchargé.
		\item Choisissez \textbf{Modify}, avancez jusqu'à l'écran «~Advanced Options~», et cochez \textbf{Add Python to environment variables}.
		\item Fermez et rouvrez votre terminal (obligatoire, sinon le changement n'est pas pris en compte).
	\end{enumerate}
	
	\subsection{Erreur~: «~pip n'est pas reconnu~»}
	
	\textbf{Solution~:} utilisez la forme complète, qui passe par Python lui-même~:
	\begin{itemize}
		\item Windows~: \texttt{python -m pip install spyder}
		\item Mac/Linux~: \texttt{python3 -m pip install spyder}
	\end{itemize}
	
	\subsection{Erreur~: «~ModuleNotFoundError: No module named '...' \,»}
	
	\textbf{Cause~:} vous essayez d'utiliser une bibliothèque (par exemple \texttt{pandas}) qui n'est pas installée.
	
	\textbf{Solution~:} installez-la avec pip, en remplaçant \texttt{nom\_module} par le nom manquant~:
	\begin{lstlisting}[language=bash]
		pip install nom_module
	\end{lstlisting}
	
	\subsection{Plusieurs versions de Python se marchent dessus}
	
	\textbf{Symptôme~:} \texttt{python --version} affiche une version différente de celle que vous venez d'installer, ou Spyder utilise une bibliothèque différente de celle que vous avez installée.
	
	\textbf{Cause~:} plusieurs installations de Python coexistent (par exemple une version système et une version Anaconda), et le PATH pointe vers la mauvaise.
	
	\textbf{Solution~:} n'installez qu'\emph{une seule} méthode (pip \emph{ou} Anaconda), pas les deux en parallèle au début. Si le mal est fait, le plus simple pour un débutant est de désinstaller Python entièrement (via le Panneau de configuration sur Windows, ou en supprimant Anaconda avec son propre désinstallateur), puis de recommencer l'installation depuis le début de ce guide, une seule fois.
	
	\subsection{Spyder ne se lance pas}
	
	\textbf{Vérifiez, dans l'ordre~:}
	\begin{itemize}
		\item Python est bien installé (\texttt{python --version} fonctionne)
		\item pip est disponible (\texttt{pip --version} fonctionne)
		\item Réinstallez Spyder proprement~: \texttt{pip install --upgrade spyder}
		\item En dernier recours, redémarrez l'ordinateur -- certaines mises à jour du PATH n'ont d'effet qu'après un redémarrage complet.
	\end{itemize}
	
	\textbf{Permission refusée sur Mac/Linux~:} si une commande d'installation échoue avec «~Permission denied~», ajoutez \texttt{sudo} devant (Linux) ou ajoutez l'option \texttt{--user}~:
	\begin{lstlisting}[language=bash]
		pip install --user spyder
	\end{lstlisting}
	
	Si rien de tout cela ne fonctionne, consultez le guide officiel~: \url{https://docs.spyder-ide.org/current/installation.html}
	
	% ============================================================
	% SECTION 7 : CONCLUSION
	% ============================================================
	\section{Conclusion}
	
	Vous avez maintenant Python et Spyder, correctement installés et vérifiés. Tout est prêt~!
	
	\subsection{La suite}
	
	\begin{itemize}
		\item Écrire des programmes Python de plus en plus complets (boucles, fonctions, fichiers)
		\item Analyser des données avec Pandas
		\item Créer des graphiques avec Matplotlib
		\item Et bientôt, du Machine Learning et du Deep Learning~!
	\end{itemize}
	
	\begin{center}
		\fbox{\parbox{0.8\textwidth}{%
				\centering
				\textbf{«~L'installation, c'est 50\% du travail. Le code, c'est le reste.~»}
				
				\smallskip
				\small Si vous bloquez, relisez ce guide. Si ça ne marche toujours pas,
				consultez la documentation officielle de Spyder.
		}}
	\end{center}
	
	\vspace{1cm}
	\begin{center}
		\LARGE \textbf{Allez, ouvrez Spyder et codez~!}
	\end{center}
	
	% ============================================================
	% ANNEXE 1 : RACCOURCIS ET COMMANDES
	% ============================================================
	\section*{Annexe~: raccourcis et commandes}
	\addcontentsline{toc}{section}{Annexe~: raccourcis et commandes}
	
	\subsection*{Raccourcis Spyder}
	
	\begin{center}
		\begin{tabular}{|l|l|}
			\hline
			\textbf{Action} & \textbf{Raccourci} \\
			\hline
			Nouveau fichier & \texttt{Ctrl+N} \\
			Ouvrir & \texttt{Ctrl+O} \\
			Enregistrer & \texttt{Ctrl+S} \\
			Exécuter tout le fichier & \texttt{F5} \\
			Exécuter la sélection/une ligne & \texttt{F9} \\
			Commenter/décommenter & \texttt{Ctrl+1} \\
			Effacer la console & \texttt{Ctrl+L} \\
			\hline
		\end{tabular}
	\end{center}
	
	\subsection*{Commandes principales}
	
	\begin{center}
		\begin{tabular}{|l|l|}
			\hline
			\textbf{Commande} & \textbf{Utilisation} \\
			\hline
			\texttt{python --version} & Vérifier la version de Python \\
			\texttt{pip --version} & Vérifier la version de pip \\
			\texttt{pip install spyder} & Installer Spyder \\
			\texttt{pip install --upgrade spyder} & Mettre à jour Spyder \\
			\texttt{pip install nom\_module} & Installer une bibliothèque \\
			\texttt{spyder} & Lancer Spyder \\
			\hline
		\end{tabular}
	\end{center}
	
	\subsection*{Petit glossaire récapitulatif}
	
	\begin{center}
		\begin{tabular}{|l|p{9.5cm}|}
			\hline
			\textbf{Terme} & \textbf{Définition courte} \\
			\hline
			Terminal & Fenêtre où l'on tape des commandes texte \\
			PATH & Liste des endroits où l'ordinateur cherche les programmes \\
			pip & Outil qui installe des bibliothèques Python \\
			IDE & Environnement de développement (ex. Spyder) \\
			Package / module & Bibliothèque de code réutilisable \\
			Script & Fichier \texttt{.py} contenant du code Python \\
			\hline
		\end{tabular}
	\end{center}
	
	% ============================================================
	% ANNEXE 2 : NOTE TECHNIQUE SUR LES IMAGES
	% ============================================================
	\section*{Note technique~: ajouter de vraies captures d'écran}
	\addcontentsline{toc}{section}{Note technique~: ajouter de vraies captures d'écran}
	
	Ce document utilise des schémas dessinés en TikZ à la place de captures d'écran, afin de rester autonome (aucun fichier externe requis pour compiler). Si vous souhaitez insérer de vraies captures d'écran officielles, procédez ainsi~:
	
	\begin{enumerate}
		\item Créez un dossier \texttt{img/} à côté de ce fichier \texttt{.tex}.
		\item Téléchargez vous-même les images souhaitées (par exemple depuis \url{https://www.python.org/downloads/} ou \url{https://docs.spyder-ide.org/}) et enregistrez-les dans \texttt{img/}, par exemple sous les noms \texttt{python-path.png}, \texttt{spyder-interface.png}.
		\item Ajoutez \texttt{\textbackslash graphicspath\{\{img/\}\}} juste après le préambule des packages.
		\item Remplacez le bloc \texttt{tikzpicture} correspondant par~: \\
		\texttt{\textbackslash capture\{python-path.png\}\{Légende\}\{label\}}
	\end{enumerate}
	
	LaTeX ne peut pas charger une image directement depuis une URL~: le fichier doit être présent sur votre disque au moment de la compilation.
	
	% ============================================================
	% REMERCIEMENTS
	% ============================================================
	\section*{Remerciements}
	\addcontentsline{toc}{section}{Remerciements}
	
	Merci à Youssouf, Soirfane, Hadidja, et surtout à vous d'avoir lu jusqu'au bout~!
	
	\begin{center}
		\textit{«~Un bon programmeur sait installer son environnement sans tout casser.~»}
		
		\smallskip
		--- Abdallah Houmadi SOIRFANE
	\end{center}
	
\end{document}